\input{preamble}

\section*{Photoelectric effect}

Let $\psi_a(\mathbf r)$ be the ground state for hydrogen.
\begin{equation*}
\psi_a(\mathbf r)=\frac{1}{\sqrt{\pi a^3}}\exp(-r/a)
\end{equation*}

Let $\psi_b(\mathbf r)$ be the wave function for a free electron.
\begin{equation*}
\psi_b(\mathbf r)=\frac{1}{\sqrt{l^3}}\exp(i\mathbf k\cdot\mathbf r)
=\frac{1}{\sqrt{l^3}}\exp(ikr\cos\theta)
\end{equation*}

Find the transition rate $R_{a\rightarrow b}$ using Fermi's golden rule
\begin{equation*}
R_{a\rightarrow b}=\frac{2\pi}{\hbar}\left|\langle\psi_b|H_1|\psi_a\rangle\right|^2\rho(E_b)
\end{equation*}

\subsubsection*{\S1}

Start by finding the perturbing Hamiltonian $H_1$.

\bigskip
We have for a particle with charge $q$
\begin{equation*}
H=\frac{1}{2m}(\mathbf p-q\mathbf A)\cdot(\mathbf p-q\mathbf A)
\end{equation*}

Expand the inner product.
\begin{equation*}
H=\frac{1}{2m}(\mathbf p^2-\mathbf p\cdot q\mathbf A-q\mathbf A\cdot\mathbf p+q^2\mathbf A^2)
\end{equation*}

Hence the first-order perturbation is
\begin{equation*}
H_1=-\frac{q}{2m}(\mathbf p\cdot\mathbf A+\mathbf A\cdot\mathbf p)
\end{equation*}

For an electron we have $q=-e$.
\begin{equation*}
H_1=\frac{e}{2m}(\mathbf p\cdot\mathbf A+\mathbf A\cdot\mathbf p)
\end{equation*}

For any wave function $\psi(\mathbf r,t)$
\begin{equation*}
(\mathbf p\cdot\mathbf A+\mathbf A\cdot\mathbf p)\psi(\mathbf r,t)
=2\mathbf A\cdot\mathbf p\psi(\mathbf r,t)-i\hbar\psi(\mathbf r,t)\nabla\cdot\mathbf A
\tag{1}
\end{equation*}

Hence for the Coulomb gauge $\nabla\cdot\mathbf A=0$ we have the identity
\begin{equation*}
\mathbf p\cdot\mathbf A+\mathbf A\cdot\mathbf p=2\mathbf A\cdot\mathbf p
\end{equation*}

Hence
\begin{equation*}
H_1=\frac{e}{m}\mathbf A\cdot\mathbf p
\end{equation*}

\subsubsection*{\S2}

Let $\mathbf E$ be the electric field
\begin{equation*}
\mathbf E=E_0\sin(\mathbf k\cdot\mathbf r-\omega t)\,\boldsymbol\epsilon
\end{equation*}

where $\boldsymbol\epsilon$ is a polarization vector such that $|\boldsymbol\epsilon|=1$
and $\mathbf k\cdot\boldsymbol\epsilon=0$.

\bigskip
Then the magnetic vector potential $\mathbf A$ is
\begin{equation*}
\mathbf A=-\frac{E_0}{\omega}\cos(\mathbf k\cdot\mathbf r-\omega t)\,\boldsymbol{\epsilon}
\end{equation*}

so that in SI units
\begin{equation*}
\mathbf E=-\frac{\partial}{\partial t}\mathbf A
\tag{2}
\end{equation*}

Hence
\begin{equation*}
H_1=\frac{e}{m}\mathbf A\cdot\mathbf p=-\frac{eE_0}{m\omega}
\cos(\mathbf k\cdot\mathbf r-\omega t)
\,\boldsymbol{\epsilon}\cdot\mathbf p
\end{equation*}

In exponential form
\begin{equation*}
H_1=
-\frac{eE_0}{2m\omega}\exp(i\mathbf k\cdot\mathbf r-i\omega t)
\,\boldsymbol{\epsilon}\cdot\mathbf p
-\frac{eE_0}{2m\omega}\exp(-i\mathbf k\cdot\mathbf r+i\omega t)
\,\boldsymbol{\epsilon}\cdot\mathbf p
\end{equation*}

Use the dipole approximation $\exp(i\mathbf k\cdot\mathbf r)\approx1$.
\begin{equation*}
H_1=
-\frac{eE_0}{2m\omega}\exp(-i\omega t)
\,\boldsymbol{\epsilon}\cdot\mathbf p
-\frac{eE_0}{2m\omega}\exp(i\omega t)
\,\boldsymbol{\epsilon}\cdot\mathbf p
\end{equation*}

Go back to cosine form.
\begin{equation*}
H_1=-\frac{eE_0}{m\omega}\cos(\omega t)\,\boldsymbol{\epsilon}\cdot\mathbf p
\end{equation*}

By the identity
\begin{equation*}
\mathbf p=\frac{im}{\hbar}(H_0\mathbf r-\mathbf rH_0)
\end{equation*}

we have
\begin{equation*}
\langle\psi_b|\mathbf p|\psi_a\rangle
=\frac{im}{\hbar}(E_b-E_a)\langle\psi_b|\mathbf r|\psi_a\rangle
\end{equation*}

Hence
\begin{equation*}
H_1=-ieE_0\frac{\omega_0}{\omega}\cos(\omega t)\,\boldsymbol\epsilon\cdot\mathbf r
\end{equation*}

where
\begin{equation*}
\omega_0=\frac{E_b-E_a}{\hbar}
\end{equation*}

For the photoelectric effect we have $\omega_0=\omega$ hence
\begin{equation*}
H_1=-ieE_0\cos(\omega t)\,\boldsymbol\epsilon\cdot\mathbf r
\end{equation*}

Rotate the lab frame so that $\boldsymbol\epsilon\cdot\mathbf r=z$.
\begin{equation*}
H_1=-ieE_0\cos(\omega t)\,z
\end{equation*}

\subsubsection*{\S3}

For perturbing Hamiltonians of the form
\begin{equation*}
H_1=V(\mathbf r)\cos(\omega t)
\end{equation*}

we have
\begin{equation*}
\langle\psi_b|H_1|\psi_a\rangle
=\frac{1}{2}\int\psi_b^*V(\mathbf r)\psi_a\,d\mathbf r
\end{equation*}

Hence for
\begin{equation*}
V(\mathbf r)=-ieE_0z
\end{equation*}

we have in polar coordinates
\begin{equation*}
\langle\psi_b|H_1|\psi_a\rangle
=-\frac{ieE_0}{2}
\frac{1}{\sqrt{\pi a^3l^3}}
\int_0^\infty\int_0^\pi\int_0^{2\pi}
\exp(ikr\cos\theta-r/a)\,zr^2\sin\theta\,dr\,d\theta\,d\phi
\end{equation*}

Integrate over $\phi$ (multiply by $2\pi$).
\begin{equation*}
\langle\psi_b|H_1|\psi_a\rangle
=-i\pi eE_0
\frac{1}{\sqrt{\pi a^3l^3}}
\int_0^\infty\int_0^\pi
\exp(ikr\cos\theta-r/a)\,zr^2\sin\theta\,dr\,d\theta
\end{equation*}

Change of variable $u=\cos\theta$.
\begin{equation*}
\langle\psi_b|H_1|\psi_a\rangle
=-i\pi eE_0
\frac{1}{\sqrt{\pi a^3l^3}}
\int_0^\infty\int_{-1}^1
\exp(ikru-r/a)\,zr^2\,dr\,du
\end{equation*}

Solve the integral over $u$.
\begin{equation*}
\langle\psi_b|H_1|\psi_a\rangle
=-i\pi eE_0
\frac{1}{\sqrt{\pi a^3l^3}}
\int_0^\infty
\frac{1}{ikr}
\left[\exp(ikr-r/a)-\exp(-ikr-r/a)\right]
zr^2\,dr
\tag{3}
\end{equation*}

Cancel $i$ and $r$.
\begin{equation*}
\langle\psi_b|H_1|\psi_a\rangle
=-\pi eE_0
\frac{1}{\sqrt{\pi a^3l^3}}\frac{1}{k}
\int_0^\infty
\left[\exp(ikr-r/a)-\exp(-ikr-r/a)\right]
zr\,dr
\end{equation*}

Substitute $z=r\cos\theta$.
This $\theta$ is different from $\theta$ in $\mathbf k\cdot\mathbf r=kr\cos\theta$.
\begin{equation*}
\langle\psi_b|H_1|\psi_a\rangle
=-\pi eE_0
\frac{1}{\sqrt{\pi a^3l^3}}\frac{\cos\theta}{k}
\int_0^\infty
\left[\exp(ikr-r/a)-\exp(-ikr-r/a)\right]
r^2\,dr
\end{equation*}

Solve the integral.
\begin{equation*}
\langle\psi_b|H_1|\psi_a\rangle
=-\pi eE_0
\frac{1}{\sqrt{\pi a^3l^3}}\frac{\cos\theta}{k}
\left[-\frac{2}{(ik-1/a)^3}+\frac{2}{(-ik-1/a)^3}\right]
\tag{4}
\end{equation*}

Rewrite as
\begin{equation*}
\langle\psi_b|H_1|\psi_a\rangle
=-\pi eE_0
\frac{1}{\sqrt{\pi a^3l^3}}\cos\theta\,
ia^4\left[\frac{16}{(a^2k^2+1)^3}-\frac{4}{(a^2k^2+1)^2}\right]
\tag{5}
\end{equation*}

\subsubsection*{\S4}

For the transition density we have
\begin{equation*}
|\langle\psi_b|H_1|\psi_a\rangle|^2=\pi e^2E_0^2
\frac{a^5}{l^3}\cos^2\theta
\left[\frac{16}{(a^2k^2+1)^3}-\frac{4}{(a^2k^2+1)^2}\right]^2
\tag{6}
\end{equation*}

Hence
\begin{align*}
R_{a\rightarrow b}&=\frac{2\pi}{\hbar}\left|\langle\psi_b|H_1|\psi_a\rangle\right|^2\rho(E_b)
\\[1ex]
&=\frac{2\pi^2}{\hbar}e^2E_0^2
\frac{a^5}{l^3}\cos^2\theta
\left[\frac{16}{(a^2k^2+1)^3}-\frac{4}{(a^2k^2+1)^2}\right]^2
\rho(E_b)
\end{align*}

The density of states for a free electron is
\begin{equation*}
\rho(E_b)
=\left(\frac{l}{2\pi}\right)^3
\frac{\sqrt{2m^3E_b}}{\hbar^3}
=\left(\frac{l}{2\pi}\right)^3
\frac{k}{a\alpha\hbar c}
\tag{7}
\end{equation*}

Hence
\begin{equation*}
R_{a\rightarrow b}
=\frac{ka^4}{4\pi\alpha\hbar^2c}
e^2E_0^2
\cos^2\theta
\left[\frac{16}{(a^2k^2+1)^3}-\frac{4}{(a^2k^2+1)^2}\right]^2
\tag{8}
\end{equation*}

\iffalse

Substitute $e^2=4\pi\varepsilon_0\alpha\hbar c$ to obtain
\begin{equation*}
R_{a\rightarrow b}
=\frac{k}{\hbar}
\varepsilon_0E_0^2a^4
\cos^2\theta
\left[\frac{16}{(a^2k^2+1)^3}-\frac{4}{(a^2k^2+1)^2}\right]^2
\end{equation*}

\fi

\subsubsection*{\S5}

Verify dimensions.
\begin{align*}
\alpha&=[1]
\\
a&=[\text{m}]
\\
k&=[\text{m}^{-1}]
\\
\omega&=[\text{s}^{-1}]
\\
e&=[\text{C}]
\\
E_0&=[\text{N}\,\text{C}^{-1}]=[\text{kg}\,\text{m}\,\text{s}^{-2}\,\text{C}^{-1}]
\\
\hbar&=[\text{J}\,\text{s}]=[\text{kg}\,\text{m}^2\,\text{s}^{-1}]
\\[1ex]
\frac{k}{a\alpha\hbar c}
&=\frac{[\text{m}^{-1}]}{[\text{m}]\,[\text{J}\,\text{s}]\,[\text{m}\,\text{s}^{-1}]}
=\frac{1}{[\text{J}]\,[\text{m}^3]}
\\[1ex]
R_{a\rightarrow b}
&=\frac{[\text{m}^{-1}]\,[\text{m}]^4}{[\text{J}\,\text{s}]^2
\,[\text{m}\,\text{s}^{-1}]}
\,[\text{N}]^2
=[\text{s}^{-1}]
\tag{9}
\end{align*}

\end{document}
